% f2015tjw_q03_flowchart.tex —— 2015 天津文科卷 Q3 程序框图（源 p1 栏 1）
%
% 源：.work/phase3-b10/source/2015天津文.pdf（p1 印 987）；实测清单 figure-inventory.md §F01
% 坐标换算：x_pic = x_pt − 175.534 ；y_pic = 542.040 − y_pt
%   （源页 pt、原点左上、y 向下；y 取反并平移，使图内画布 y 向上；
%     画布范围 x' 0.000–95.833、y' 0.147–204.730，单位 pt）
% 线宽：0.398 pt（源 `0.398 w` 实测；全图实线，p1 无虚线）
% 箭头：TikZ 标准 `>=stealth`（既定重绘约定）。
%   源箭尖为 8 点填充多边形（实测 2.88–3.36 宽 × 3.6–4.0 长），其内容流提取坐标与墨迹不符
%   （清单 §S-1：沿线采样仅 15.42% 落墨）。本文件**不采用提取值**，箭尖落点一律由本次
%   600 dpi 墨迹独立实测（口径：沿轴扫描，取墨迹最远端；见 .work/phase3-b10/probe/arrowtips 系）：
%     C1 起 (216.223, 354.347) 尖 368.16 ｜ C2 起 (216.223, 385.120) 尖 398.88
%     C3 起 (216.223, 414.491) 尖 428.28 ｜ C4 起 (216.223, 444.113) 尖 457.92
%     C5 起 (216.223, 479.071) 尖 492.84 ｜ C6 起 (216.223, 510.504) 尖 524.28
%     否支折线 3 段，尖 x = 216.96（源提取值 223.396 系偏移值，已弃用）
%   （六枚主链箭尖与目标框边的墨迹间隙均 0.36 pt；C4 与菱形上顶点同）
% 存疑照源复刻：否支出口 N1 起点 x=257.194 与菱形右顶点 256.913 差 0.281 pt；
%   否支落点在 B2/B3 之间的竖段上（y=392.206）；标签「是」在菱形下方、「否」在菱形右侧 —— 均按源。
\begin{tikzpicture}[x=1pt, y=1pt, line width=0.398pt, >=stealth, every node/.style={font=\fontsize{10.5pt}{12.6pt}\selectfont,inner sep=0pt}]

  % ---- 七个框 ----
  \draw[rounded corners=3.985pt] (26.690,204.730) rectangle (54.689,187.892);   % B1 开始
  \draw (6.776,173.320) rectangle (74.603,157.119);                             % B2 S=10,i=0
  \draw (17.530,142.547) rectangle (63.848,127.748);                            % B3 i=i+1
  \draw (14.920,113.176) rectangle (66.458,98.127);                             % B4 S=S-i
  \draw (81.379,73.362) -- (40.689,83.473) -- (0.000,73.362)
        -- (40.689,63.251) -- cycle;                                            % B5 菱形 S<=1?
  \draw (15.281,31.736) -- (21.417,48.595) -- (66.092,48.595)
        -- (59.957,31.736) -- cycle;                                            % B6 输出 S（平行四边形）
  \draw[rounded corners=3.985pt] (26.690,17.163) rectangle (54.689,0.147);      % B7 结束

  % ---- 主链 6 段（箭尖落点见文件头实测值） ----
  \draw[->] (40.689,187.693) -- (40.689,173.880);   % 354.347 -> 尖 368.16
  \draw[->] (40.689,156.920) -- (40.689,143.160);   % 385.120 -> 尖 398.88
  \draw[->] (40.689,127.549) -- (40.689,113.760);   % 414.491 -> 尖 428.28
  \draw[->] (40.689, 97.927) -- (40.689, 84.120);   % 444.113 -> 尖 457.92
  \draw[->] (40.689, 62.969) -- (40.689, 49.200);   % 479.071 -> 尖 492.84
  \draw[->] (40.689, 31.536) -- (40.689, 17.760);   % 510.504 -> 尖 524.28

  % ---- 「否」支：3 折段回环，落于 B2 与 B3 之间竖段 ----
  \draw[->] (81.660,73.362) -- (95.833,73.362) -- (95.833,149.834) -- (41.426,149.834);

  % ---- 标签（源 pdftotext -bbox 字形盒左上角；inner sep=0） ----
  \node[anchor=north west,inner sep=0] at (30.186,202.140) {开始};
  \node[anchor=north west,inner sep=0] at (10.276,169.930) {$S=10,\,i=0$};
  \node[anchor=north west,inner sep=0] at (21.026,139.410) {$i=i+1$};
  \node[anchor=north west,inner sep=0] at (18.416,109.790) {$S=S-i$};
  \node[anchor=north west,inner sep=0] at (25.066, 80.350) {$S\leqslant 1\,?$};
  \node[anchor=north west,inner sep=0] at (44.386, 61.820) {是};
  \node[anchor=north west,inner sep=0] at (76.406, 70.890) {否};
  \node[anchor=north west,inner sep=0] at (24.916, 46.030) {输出 $S$};
  \node[anchor=north west,inner sep=0] at (30.186, 14.510) {结束};

\end{tikzpicture}
